\section{Van de Vusse System Identification}

This report adds the Van de Vusse CSTR as the third case-study seed in the repository, but limits the implementation to the nonlinear plant and system-identification layer. No MPC or RL workflow is introduced in this phase.

\subsection{Benchmark Role}

The benchmark is attractive for later assisted-MPC work because it combines:
\begin{itemize}
    \item fast nonlinear dynamics,
    \item thermal and concentration coupling,
    \item disturbances through the inlet concentration,
    \item nonminimum-phase behavior in the relevant operating region, and
    \item a compact two-input, two-output control surface.
\end{itemize}

The active plant outputs for this repo are chosen as
\[
y = \begin{bmatrix} c_B \\ T \end{bmatrix},
\]
which provides one quality variable and one thermal variable for the later MPC layer.

\subsection{Reaction System and Nonlinear Plant}

The Van de Vusse reaction network is
\[
A \rightarrow B,\qquad
B \rightarrow C,\qquad
2A \rightarrow D.
\]

The nonlinear plant is implemented with the Kelvin-form state vector
\[
x = \begin{bmatrix} c_A & c_B & T & T_K \end{bmatrix}^\top
\]
and manipulated inputs
\[
u = \begin{bmatrix} F & Q_K \end{bmatrix}^\top.
\]

The model is
\[
\frac{dc_A}{dt} = F(c_{A0} - c_A) - k_1(T)c_A - k_3(T)c_A^2
\]
\[
\frac{dc_B}{dt} = -Fc_B + k_1(T)c_A - k_2(T)c_B
\]
\[
\frac{dT}{dt} = F(T_{in} - T)
- \frac{k_1(T)c_A \Delta H_{AB} + k_2(T)c_B \Delta H_{BC} + k_3(T)c_A^2 \Delta H_{AD}}{\rho C_p}
+ \frac{k_w A_R}{\rho C_p V_R}(T_K - T)
\]
\[
\frac{dT_K}{dt} = \frac{Q_K + k_w A_R(T - T_K)}{m_K C_{pK}}
\]
with Arrhenius-rate definitions
\[
k_i(T) = k_{i0}\exp\!\left(\frac{E_{ai}}{T}\right).
\]

The ODE structure and nominal parameter table follow the PINN-oriented modeling paper, while the benchmark operating point and constraints follow the Klatt/Engell control benchmark description.

\subsection{Benchmark Operating Point and Constraints}

The local identification workflow is centered on the benchmark operating inputs
\[
F_s = 18.83\ \mathrm{h}^{-1}, \qquad
Q_{K,s} = -4495.7\ \mathrm{kJ/h}.
\]

The benchmark state seed used for the steady-state solve is
\[
c_{A,s} = 1.235,\quad
c_{B,s} = 0.9,\quad
T_s = 407.29\ \mathrm{K},\quad
T_{K,s} = 402.10\ \mathrm{K}.
\]

The implemented plant still solves its own steady state from the nonlinear PINN-parameterized model and the chosen input pair. The benchmark point is therefore used as a physically meaningful seed, not as a hard-coded plant state.

The adopted operating constraints are:
\begin{itemize}
    \item \(5 \leq F \leq 35\ \mathrm{h}^{-1}\)
    \item \(-8500 \leq Q_K \leq 0\ \mathrm{kJ/h}\)
    \item \(4.5 \leq c_{A0} \leq 5.7\ \mathrm{mol/L}\)
    \item \(0.7 \leq c_B \leq 0.95\ \mathrm{mol/L}\)
\end{itemize}

\subsection{Step-Test Design}

The control benchmark does not prescribe system-identification steps directly, so the system-ID notebook uses local signed steps around the operating point while remaining well inside the benchmark bounds.

The canonical sample time is
\[
\Delta t = 0.01\ \mathrm{h}
\]
with:
\begin{itemize}
    \item initial hold \(= 0.2\ \mathrm{h}\)
    \item stepped hold \(= 1.2\ \mathrm{h}\)
\end{itemize}

The eight fitting tests are:
\begin{itemize}
    \item \(F \pm 2.0\)
    \item \(F \pm 4.0\)
    \item \(Q_K \pm 500\)
    \item \(Q_K \pm 1000\)
\end{itemize}

One extra validation case is also saved:
\[
\Delta F = +2.0,\qquad \Delta Q_K = -500.
\]

These steps are deliberately local. The goal is a clean nominal linear model for future offset-free MPC, not a global nonlinear approximation.

\subsection{FOPDT Identification}

Each input-output channel is identified with a delayed first-order model
\[
G_{ij}(s) = \frac{K_{ij}}{\tau_{ij}s + 1} e^{-\theta_{ij}s}.
\]

The implementation keeps the current repository spirit by using a 28/63-style heuristic, but it is fully automatic and non-interactive. For each saved step test the notebook estimates:
\begin{itemize}
    \item channel gain \(K_{ij}\),
    \item time constant \(\tau_{ij}\),
    \item delay \(\theta_{ij}\).
\end{itemize}

The per-input channel estimates are aggregated across the signed small and medium steps using the median, which keeps the first-pass workflow simple and robust to one-off nonlinear asymmetries.

\subsection{Python Delay Absorption}

The new Van de Vusse workflow does not call MATLAB \texttt{absorbDelay}. Instead, each FOPDT channel is converted to a discrete-time realization directly in Python.

For the non-delayed first-order part:
\[
a_{ij} = e^{-T_s/\tau_{ij}}, \qquad
b_{ij} = K_{ij}(1-a_{ij}).
\]

This yields the scalar dynamic state
\[
x_{ij}[k+1] = a_{ij} x_{ij}[k] + b_{ij} v_{ij}[k].
\]

The delay is quantized to an integer number of samples
\[
m_{ij} = \mathrm{round}(\theta_{ij}/T_s)
\]
and absorbed with a shift register:
\[
d_{ij,1}[k+1] = u_j[k], \qquad
d_{ij,\ell+1}[k+1] = d_{ij,\ell}[k].
\]

If \(m_{ij}=0\), then \(v_{ij}[k]=u_j[k]\). Otherwise, \(v_{ij}[k]\) is the last delay-chain state. The global MIMO realization stacks one dynamic state plus the needed delay states for each channel and sums the channel contributions into each measured output.

This keeps the Van de Vusse system-identification path fully inspectable and removes any new MATLAB dependency for the new case study.

\subsection{Saved Artifacts}

The canonical notebook saves its outputs under:
\begin{itemize}
    \item \texttt{VanDeVusse/Data}
    \item \texttt{VanDeVusse/Results}
\end{itemize}

The main saved artifacts are:
\begin{itemize}
    \item nine step-test CSV files,
    \item \texttt{system\_dict.pickle} with identified \(A, B, C, D\),
    \item \texttt{scaling\_factor.pickle},
    \item \texttt{min\_max\_states.pickle},
    \item metadata pickle/json describing the fitted channels and delay steps,
    \item comparison plots between the nonlinear plant and the identified model.
\end{itemize}

\subsection{Next Step}

This report intentionally stops at the plant-plus-identification stage. The next phase can use the saved Van de Vusse matrices and scaling artifacts to build a canonical offset-free MPC notebook and then RL-assisted supervisory layers, following the same overall progression already used for polymer and distillation.
